% Generated draft: requires mathematical and source review.
\documentclass[11pt]{article}
\usepackage[margin=25mm]{geometry}
\usepackage{fontspec}
\setmainfont{Latin Modern Roman}
\usepackage{amsmath,amssymb,mathtools,unicode-math}
\setmathfont{Latin Modern Math}
\usepackage{enumitem,needspace,xurl,hyperref}
\setlength{\parindent}{0pt}
\setlength{\parskip}{5pt}
\setlength{\emergencystretch}{3em}
\newcommand{\N}{\mathbb{N}}
\newcommand{\Z}{\mathbb{Z}}
\newcommand{\Q}{\mathbb{Q}}
\newcommand{\R}{\mathbb{R}}
\newcommand{\C}{\mathbb{C}}
\newcommand{\F}{\mathbb{F}}
\newcommand{\A}{\mathbb{A}}
\newcommand{\PP}{\mathbb{P}}
\newcommand{\E}{\mathbb{E}}
\newcommand{\eps}{\varepsilon}
\newcommand{\dd}{\,\mathrm{d}}
\begin{document}
\section*{UnsolvedMath ID: 30006518}
\textbf{Source problem code:} catalogue entry
\subsection{Definitions and mathematical statement}
For a first-order structure \(M\), let \(\operatorname{Th}(M)\) denote its complete theory. Say that \(M\) is \(\omega\)-categorical if \(\operatorname{Th}(M)\) has, up to isomorphism, exactly one countable model. Let \(\operatorname{Aut}(M)\) be its automorphism group, with composition as the group operation. Define its elementary self-embedding monoid by
\[
\operatorname{EEmb}(M)=\{f:M\to M : f\text{ is injective and elementary}\}.
\]
Here elementary means that, for every first-order formula \(\varphi(x_1,\ldots,x_n)\) in the language of \(M\) and every tuple \(\bar a\in M^n\),
\[
M\models\varphi(\bar a)\quad\Longleftrightarrow\quad M\models\varphi(f(\bar a)).
\]
The monoid operation is composition and its identity is \(\operatorname{id}_M\). The topology of pointwise convergence on either invariant has basic neighborhoods
\[
U(f,F)=\{g:g\vert_F=f\vert_F\},\qquad F\subseteq M\text{ finite},
\]
where \(g\) ranges over the relevant group or monoid.

The supplied background suggests reconstruction up to bi-interpretability. In the parameter-free reading, \(M\) and \(N\) are bi-interpretable if each admits a first-order interpretation in the other, without parameters, and the two composite interpretations are definably isomorphic to the respective original structures. Write \(M\simeq_{\mathrm{bi}}N\) for this relation.

\textbf{Reconstruction question.} Determine additional hypotheses and a comparison class \(\mathcal C\) of \(\omega\)-categorical structures under which either of the following reconstruction properties holds, considered separately:
\[
\begin{aligned}
&\forall M,N\in\mathcal C,\qquad
\operatorname{Aut}(M)\cong_{\mathrm{grp}}\operatorname{Aut}(N)
\ \Longrightarrow\ M\simeq_{\mathrm{bi}}N;\\
&\forall M,N\in\mathcal C,\qquad
\operatorname{EEmb}(M)\cong_{\mathrm{mon}}\operatorname{EEmb}(N)
\ \Longrightarrow\ M\simeq_{\mathrm{bi}}N.
\end{aligned}
\]
The isomorphisms in these implications are purely algebraic; continuity for the topology of pointwise convergence is not an assumption.

\textbf{Unresolved scope.} These implications give a precise scheme for the supplied question, not a conjecture that either implication holds for all \(\omega\)-categorical structures. The record does not specify the additional hypotheses, the comparison class, whether those hypotheses must hold for both structures, or the exact reconstruction equivalence intended in the shortened statement. The displayed scheme uses the bi-interpretability target suggested by the original background and explicitly chooses its parameter-free version. The customary setting of countable structures in countable first-order languages also requires confirmation: neither the cardinality of \(M\) nor the language restrictions are explicit in the shortened statement. If that setting is intended, it must be included in the specification of \(\mathcal C\). Whether finite structures are included is likewise unspecified.
\subsection{Short English statement}
Determine conditions under which the abstract automorphism group, or separately the abstract monoid of elementary self-embeddings, determines an \(\omega\)-categorical structure up to bi-interpretability, without supplying the topology of pointwise convergence. The comparison class, language and cardinality conventions, and parameter convention require confirmation. This is a statement-only draft requiring human mathematical review.
\subsection{Sources}
\begin{itemize}
\item Supplied local UnsolvedMath registry record, permanent ID \(30006518\), source problem code \texttt{catalogue entry}, title \emph{Abstract-Group Reconstruction of Countably Categorical Structures}, in the collection \emph{Oberwolfach Reports Open Problems}. No proposer is specified.
\item Supplied report citation: \emph{Homogeneous Structures: Model Theory meets Universal Algebra} (2026), Oberwolfach Reports, pp.\ 3063--3124, DOI \texttt{10.4171/owr/2025/58}: \url{https://doi.org/10.4171/owr/2025/58}. The extracted original text includes the heading \emph{Notes on a reconstruction theorem of Lascar}, the name Mira Tartarotti, and page 3111. This does not by itself identify the proposer. The registry gives proposed year 2025, while its report citation gives 2026.
\item The supplied original excerpt attributes topological reconstruction up to bi-interpretability to a result of Coquand reported by Ahlbrandt and Ziegler, reference [1], and independently proved by Makkai, reference [5]. Full bibliographic entries for [1] and [5] are absent from the supplied record. Its extracted sentence about purely algebraic reconstruction is truncated.
\item Supplied literature source: Bodirsky--Pinsker, \emph{Aleph-zero-categorical structures}, \url{https://arxiv.org/abs/0907.2925}.
\item Supplied literature source: Pech--Pech, \emph{Reconstructing topology of elementary self-embedding monoids}, \url{https://doi.org/10.1007/S11225-017-9756-6}.
\item Supplied literature source: \emph{Oberwolfach Report 2025/58}, \url{https://doi.org/10.4171/OWR/2025/58}.
\item Literature assessment dated exactly \texttt{2026-08-22}: the supplied curation labels the problem partially solved and states that topological reconstruction is classical, that purely algebraic reconstruction is established under several additional hypotheses, and that no general topology-free reconstruction theorem is known. This is supplied dated curation, not an independently verified literature assessment; the source explicitly requests independent verification of time-sensitive status and cited claims.
\item The registry labels the statement \texttt{corrected\_verified} and reports that the official Oberwolfach report was checked, the question isolated, and notation repaired. These are supplied source-verification claims, not verification performed for this draft; no separate verification date is supplied.
\end{itemize}
\end{document}
